2.1-4 \\
Consider the searching problem: \\
\textbf{Input}: A sequence of \textit{n} numbers $<a_1, a_2, ..., a_n>$ stored in array A[1:n] and a value x. \\
\textbf{Output}: An index \textit{i} that \textit{x} equals A[i] or the special value \textit{NIL} if \textit{x} doesn't appear in A. \\

Write pseudocode for linear search, which scans through the array from begin-ning to end, looking for 
x. Using a loop invariant, prove that your algorithm is 
correct. Make sure that your loop invariant fulûlls the three necessary properties.

\begin{verbatim}
# give the index of the x target or else -1 if not found
def find_val(x: int, array: list[int]) -> int:
    n: int = len(array)
    res: int = -1 

    for i in range(n):
        if (array[i] == x):
            res = i
            break
            
    return res
            
array: list[int] = [1,2,3,4,5,6,7]
index: int = find_val(100, array)

print(f"Index: {index}")
\end{verbatim}

Note: Loop invariant properites include:
\begin{itemize}
    \item \textbf{Initialization}: The invariant is true prior to the first iteration of the loop.
    
    \item \textbf{Maintenance}: If the invariant is true before an iteration of the loop, it remains true before the next iteration.
    
    \item \textbf{Termination}: The loop eventually terminates. When it does, the invariant—together with the reason the loop terminated—gives a useful property that helps prove the algorithm is correct.
\end{itemize}

\begin{proof}
\textbf{
    Initialization:} Before the first iteration, the loop variable $i$ is set to $0$.
At this point, the subarray $A[0:i]$ contains no elements. Therefore, no index $j < i$
exists such that $A[j] = x$. The loop invariant holds because we have not yet examined
any elements, so it is correct that $x$ has not been found in the subarray. Thus we return
Nil (-1) which is a valid solution in this scenerio. 

\textbf{Maintenance}: 

\end{proof}

2.1-5 \\
Consider the problem of adding two n-bit binary integers a and b, stored in two n-element arrays
A[0:n-1] and B[0:n-1], where each element is either 0 or 1, a = $\sum_{i= 0}^{n-1}$A[i]*$2^i$, and b = $\sum_{i=0}^{n-1}B[i]*2^i$.
The sum c = a + b of the two integers should be stored in binary from in an (n + 1)-element array C[0:n], where c = $\sum_{i=0}^{n}C[i]*2^i$.
Write a procedure ADD-BINARY-INTEGERS that takes as input arrays A and B, along with the length n, and returns array C holding the sum.

